\documentclass[a4paper, 12pt]{article}

% include required packages
\usepackage[a4paper, margin=2cm]{geometry} % Set page margins
\usepackage[round,authoryear]{natbib}      % For citations
\usepackage{url}                           % For formatting URLs
\usepackage{indentfirst}                   % Indent first paragraph of sections
\usepackage{booktabs}                      % For better table formatting
\usepackage{float}                         % For [H] table placement
\usepackage{enumitem}                      % For customising lists
\usepackage{tikz}                          % For drawing diagrams
\usetikzlibrary{positioning}               % For relative positioning

% set paragraph and list formatting
\setlength{\parindent}{1.5em}
\setlength{\parskip}{0pt}
\setlist[itemize]{noitemsep,topsep=0pt,parsep=0pt,partopsep=0pt,after=\vspace{0.5em}}


% set title information
\title{A Critical Code Review of the Web Witchcraft and Wizardry Project}
\author{Edmund Mulligan}
\date{\today}

% document start
\begin{document}

% set bibliography style and page numbering
\bibliographystyle{plainnat}
\pagenumbering{arabic}

\maketitle

\begin{abstract}
This essay presents a critical code review of the Web Witchcraft and Wizardry project, an interactive web-based learning platform. The review examines the project's modular architecture, focusing on JavaScript implementation including encrypted local storage, DOM manipulation, and defensive programming practices. The analysis demonstrates adherence to web development best practices while reflecting on the iterative development process and testing methodologies employed.
\end{abstract}

\tableofcontents
\newpage

\section{Introduction}
This essay reviews the codebase of the Web Witchcraft and Wizardry project, a web-based interactive learning platform designed to teach users about web development concepts through a magical theme. The review covers the architecture and design of the project, as well as an analysis of the HTML, CSS, and JavaScript code, with an emphasis on the latter. Reflections on the development process and potential areas for improvement are also discussed. The code is available for review on GitHub at 

\url{https://github.com/Birkbeck2/puffin-web-project-edmundmulligan}. The deployed code can be viewed in GitHub Pages or at 

\url{https://web.witchcraft-and-wizardry.school/}. Results of automated tests (along with design documentation) can be viewed at 

\url{https://web.witchcraft-and-wizardry.school/diagnostics}.

\section{Architecture and Design}
The design of the appearance and layout of the website was documented in the previous essay \citep{mulligan2026a}, so will not be covered here, but the code has also been designed so as to be maintainable and extensible, with a modular structure that separates concerns \citep{dijkstra1982role}. This is a front-end focused project, so the codebase is primarily composed of HTML, CSS, and JavaScript files, with a clear separation between structure, presentation, and behaviour \citep{keith2010domscripting}. Following DevOps principles \citep{allman2018devops}, the project utilises a build process to allow automated testing, as well as a version control system to manage changes and collaborate with other developers (should there be any in the future).

Several constraints have been imposed on the project, partly to comply with the academic requirements of this course, but also to keep the project focused and manageable.
\begin{itemize}
    \item Use only vanilla JavaScript, HTML, and CSS, without any external libraries or frameworks.
    \item No use of generative AI or Large Language Models (LLMs) in the development process.
    \item No server side code or databases, as the project is focused on front-end development and client-side interactivity.
    \item Ensure that the website is fully accessible and responsive, adhering to best practices in web development.
    \item The project is developed by a single individual (myself).
\end{itemize}

\section{HTML and CSS}
Since the focus of this assignment is Javascript, the HTML and CSS code will not be analysed in depth, but some general observations can be made. The HTML code is well-structured and semantic, with appropriate use of headings, lists, and other elements to create a clear and logical document structure. The CSS code is organised into separate files for layout and components, with consistent naming conventions and a clear separation of concerns. The use of CSS variables allows for easy theming and customization of the website's appearance. The CSS has been completely refactored since the last assignment, to make it more maintainable and efficient. This includes keeping all global variables in definition files, using layers to control the cascade and specificity of styles, as well as removing redundant and unused styles. Nested selectors have been used where possible to improve readability. At most three CSS files need to be loaded on any page (globals.css, main.css and a page specific CSS file if necessary), which makes it less likely a specific CSS file will be accidentally missed when developing new pages. The CSS code has been heavily influenced by tutorials from Kevin Powell \citep{powell_youtube} and his new paid course \citep{powell2026css}, which have provided valuable insights into modern CSS techniques and best practices.

While most html files are reasonably well structured, students/lesson-00.html is an exception, as it is so large. As an HTML file is is unmaintainable and has been  generated from a template using moustache. There is no functionality in this web page that is not better illustrated in other more reasonably sized pages, so it is requested to be excluded from the assessment of this project. It is included since it is necessary for students who want to use the web site.

\section{Javascript}
Arguably, there is rather more JavaScript in this project than is required for academic purposes, but this is because the project is designed to be a fully functional and interactive web application, rather than just a demonstration of specific concepts. The JavaScript code is organised into separate files for different components and features. The code follows best practices for readability and maintainability \citep{McConnell2004}, with consistent naming conventions, clear comments, and a focus on simplicity and clarity. 

There is insufficient space in this essay to discuss all the JavaScript code in detail so this section highlights some of the more interesting and complex aspects. See Appendix \ref{appendix:js} for a sample of examples of JavaScript code that demonstrate the use of various concepts and techniques required by this assessment.

\subsection{localStorage.js}
Arguably this is the most complex Javascript file in the codebase. It was written after a code review using CodeQL \citep{codeql2024} identified that using unencrypted local storage was a security risk. Arguably, no sensitive data is being stored so this is a false positive, but best practice is to eliminate security flaws before they become a problem \citep{owasp2010secure}. It uses AES-GCM 256-bit encryption \citep{nist2001aes} to encrypt data before storing it in localStorage, and decrypts it when retrieving. The encryption key is derived from a passphrase using PBKDF2 with a random salt and the encrypted data is stored as a JSON string containing the ciphertext and salt.

\subsection{populateLineNumbers.js}
This code is used to display javascript code snippets on the website with line numbers. It works by selecting all elements with the class 'code-snippet', splitting their text content into lines, and then creating a new HTML structure that includes both the line numbers and the code lines. The line numbers are generated dynamically based on the number of lines in the code snippet, and the original code is preserved in a separate element for styling purposes. This allows for separation of concerns between the structure of the code and its presentation, as well as making it easier to maintain and update the code snippets in the future. It also allows for the code snippets to be easily copied without including the line numbers, which can be a common issue when displaying code on websites.

\subsection{injectGlossaryPopover.js}
This code fetches glossary data from glossary.html and injects it into the page as popovers that appear when hovering over glossary terms. It works by selecting all elements with the class 'glossary-term', retrieving the corresponding definition from the glossary data, and then creating a popover element that is displayed on hover. To prevent injection attacks, the code sanitises the glossary definitions by escaping any HTML tags before inserting them into the popover. This ensures that any potentially malicious code in the glossary definitions cannot be executed when the popover is displayed. Rendering of the popover is handled by CSS, which, as usual, allows for separation of structure and presentation.

\subsection{utils.js, queryParameters.js and debug.js}
These scripts are not directly related to any specific features of the website, but provide utility functions and debugging tools that are used throughout the codebase. The utils.js file contains a variety of helper functions for tasks such as DOM manipulation, event handling, and data formatting. The queryParameters.js file provides a simple interface for working with URL query parameters, allowing for easy retrieval and manipulation of parameters in the URL. The debug.js file contains functions for logging and displaying debug information, which was essential during development and testing.

\section{Reflections}

The development process was highly iterative, with continuous refinement driven by automated testing and manual code review. Given the scope of the project, and the fact that it was developed by a single person, this was the only feasible approach and drew on my considerable experience of developing in a DevOps environment. The website is at the limit of what can be achieved without server side processing. In particular, using local storage to persist data is inferior to using a database for user experience as using a different computer will lose all saved data. Although not the focus of this assignment, I am particularly please with how the CSS refactoring turned out. The CSS submitted last term was functional but messy, reflecting the fact that I was learning it while developing the project The refactored version is much cleaner, particularly with the use of CSS variables in the definitions folder and layers. The Javascript code is reasonably well organised but could be improved. For example localStorage.js could and should, with more time, be refactored into classes to handle encryption and storage separately.

\section{Conclusion}

This code review demonstrates that the Web Witchcraft and Wizardry project achieves its academic objectives through well-structured, maintainable code adhering to web development best practices. It is at a level of maturity where it can be given to students, mentors and teachers as a working prototype to solicit feedback before being released as a final product. The deliberate architectural choices (modular JavaScript organization, semantic HTML, modern CSS and separation of concerns) create a robust foundation for future enhancement. The comprehensive testing regime validates accessibility and standards compliance. The codebase exhibits professional-grade practices including encryption for data storage and defensive programming against injection attacks. Further work is required to write the lessons and this will be the focus of the next phase of development, along with user and expert testing.

% Include word count N.B. Run wordcount.sh to generate wordcount.txt before compiling
\section*{Word count}
Word count: \input{essay-3-wordcount.txt} words (excluding references)
\\
\bibliography{EdmundLibrary}

% include appendices
\newpage
\appendix
\section*{Appendices}
\section{Javascript Assessment Criteria}\label{appendix:portraits}

Here is an index of where different Javascript concepts are used in the codebase to allow assessors to easily find examples of the required concepts. This is not an exhaustive list of all Javascript concepts used in the codebase, but rather a guide to where the key features can be found.

\subsection{Variables and Functions}

\subsection{Objects and Arrays}

\subsection{Conditional Logic and Loops}

\subsection{DOM Manipulation}

\subsection{Interactivity}

\subsection{Data Retrieval and Storage}

\end{document}
